\documentclass[stu,12pt,floatsintext]{apa7}

% Configuración de idioma (Babel Spanish)
\usepackage[spanish]{babel}
\usepackage{csquotes}
% Nota: Si aparece error de bibliografía, recuerda ejecutar Biber después de compilar con pdfLaTeX
\usepackage[style=apa,sortcites=true,sorting=nyt,backend=biber]{biblatex}
\DeclareLanguageMapping{spanish}{spanish-apa}
\addbibresource{bibliografia.bib} 

\usepackage[T1]{fontenc}
\usepackage{mathptmx} % Fuente Times New Roman [3]

% Información de la portada
\title{Reporte de Observación de Nubosidad: Proyecto de Ciencia Ciudadana GLOBE}
\shorttitle{Reporte de Nubosidad}

% Corrección: En modo 'stu', listar autores sin usar el comando \and
\author{Daniel Soto Villada, Jean Lucca Buritica Cuervo, Camilo Higuita y Rubén Páez}

\affiliation{Universidad de Antioquia}
\course{Astronomía Práctica}
\professor{Nombre del Profesor}
\duedate{1 de marzo de 2026}

\abstract{Este reporte documenta las observaciones sistemáticas de la cobertura y tipo de nubes realizadas mediante la aplicación GLOBE Observer. El objetivo es comprender la variabilidad temporal de la nubosidad y su impacto en la planeación de observaciones astronómicas [4], [5].}

\begin{document}
\maketitle

\section{Introducción}
La fracción de cielo despejado es un factor crucial en la planeación de observaciones astronómicas [4]. Este proyecto busca relacionar observaciones locales con procesos meteorológicos y validar datos satelitales a través del programa GLOBE de la NASA [4], [5].

\section{Método}
\subsection{Materiales}
Se utilizó la aplicación \textit{GLOBE Observer} para el registro de datos [5].
\subsection{Procedimiento}
Se realizaron observaciones diarias registrando ubicación, fecha, hora, cobertura y tipo de nubes [6].

\section{Resultados}
% Espacio para incluir tablas y figuras de la aplicación GLOBE.

\section{Análisis y Discusión}
% Instrucciones de las fuentes para completar esta sección [6], [7], [8]:
% - Identificar nubes más frecuentes y patrones observados.
% - Relacionar hora del día con cantidad de nubosidad.
% - Discutir dificultades de estimación (subjetividad, visibilidad).
% - Comparar con el clima típico de la región en esta época.

\printbibliography
\end{document}

